% =============================================================
%  Insulation Coordination — Part II (Module 08)
% =============================================================
\begin{frame}{Lightning stroke}
  \begin{columns}[c,onlytextwidth]
    \begin{column}{0.48\textwidth}
      \centering
      \topoimg[width=\linewidth,height=0.66\textheight,keepaspectratio]{\icfig{m08_image4}}
    \end{column}
    \begin{column}{0.48\textwidth}
      \centering
      \topoimg[width=\linewidth,height=0.66\textheight,keepaspectratio]{\icfig{m08_image5}}
    \end{column}
  \end{columns}
  \vspace{0.3em}
  {\color{atGray}\footnotesize IEEE Std 998-1996.}
\end{frame}

\begin{frame}{Lightning stroke --- travelling wave}
  \begin{columns}[c,onlytextwidth]
    \begin{column}{0.56\textwidth}
      \begin{itemize}
        \setlength{\itemsep}{0.4em}
        \item Travelling-wave propagation is not comparable to the current
              flow at nominal frequency.
        \item On an open circuit no current flows, but the travelling wave
              can still enter the conductor.
        \item When the wave hits the open circuit (e.g. a transformer) the
              surge is reflected: the voltage keeps its polarity but the
              amplitude doubles.
        \item The increased amplitude at transformers must be considered in
              the equipment design.
      \end{itemize}
      \vspace{0.2em}
      {\color{atGray}\footnotesize Westinghouse, ``Electrical Transmission
      and Distribution Reference Book'', 1961.}
    \end{column}
    \begin{column}{0.40\textwidth}
      \centering
      \topoimg[width=\linewidth,height=0.66\textheight,keepaspectratio]{\icfig{m08_image6}}
    \end{column}
  \end{columns}
\end{frame}

\begin{frame}{Overvoltages in AC grids}
  \centering
  \topoimg[width=\linewidth,height=0.72\textheight,keepaspectratio]{\icfig{m08_image7}}\\[0.3em]
  {\color{atGray}\footnotesize ``Metal-Oxide Surge Arresters in High-Voltage
  Power Systems'', Volker Hinrichsen.}
\end{frame}

\begin{frame}{Arrester design}
  \begin{columns}[c,onlytextwidth]
    \begin{column}{0.56\textwidth}
      \begin{itemize}
        \setlength{\itemsep}{0.4em}
        \item Arresters used today are metal-oxide (MO) arresters without gaps.
        \item Made from MO resistor disks connected in series.
        \item The number of series-connected resistors defines the arrester
              characteristics.
        \item MO resistors have an extremely non-linear U--I characteristic.
        \item Resistors are placed in a housing for protection against
              environmental conditions.
      \end{itemize}
    \end{column}
    \begin{column}{0.40\textwidth}
      \centering
      \topoimg[width=\linewidth,height=0.68\textheight,keepaspectratio]{\icfig{m08_image8}}
    \end{column}
  \end{columns}
\end{frame}

\begin{frame}{U--I arrester characteristic}
  \begin{columns}[c,onlytextwidth]
    \begin{column}{0.54\textwidth}
      \begin{itemize}
        \setlength{\itemsep}{0.35em}
        \item $U_S$: maximum system voltage (phase-phase AC).
        \item $U_c$: maximum continuous operating voltage (MCOV).
        \item For DC applications the crest value of continuous operating
              voltage (CCOV) is used instead of MCOV.
        \item Lightning impulse current ($I_n$): 8/20\,\textmu s
              (fast-front).
        \item Switching impulse current: 30/60\,\textmu s (slow-front).
      \end{itemize}
    \end{column}
    \begin{column}{0.42\textwidth}
      \centering
      \topoimg[width=\linewidth,height=0.68\textheight,keepaspectratio]{\icfig{m08_image9}}
    \end{column}
  \end{columns}
\end{frame}

\begin{frame}{Arrester protection capability --- connection}
  \begin{columns}[c,onlytextwidth]
    \begin{column}{0.56\textwidth}
      \begin{itemize}
        \setlength{\itemsep}{0.4em}
        \item Protection capability is based on providing a low-impedance
              path towards earth.
        \item That impedance is defined by the arrester itself and by its
              connection to the overhead line and to earth.
        \item Either accurate calculations or safety margins are applied to
              account for those effects.
      \end{itemize}
    \end{column}
    \begin{column}{0.40\textwidth}
      \centering
      \topoimg[width=\linewidth,height=0.66\textheight,keepaspectratio]{\icfig{m08_image10}}
    \end{column}
  \end{columns}
\end{frame}

\begin{frame}{Arrester protection capability --- travelling wave}
  \begin{columns}[c,onlytextwidth]
    \begin{column}{0.56\textwidth}
      \begin{itemize}
        \setlength{\itemsep}{0.4em}
        \item The surge current will enter a ``dead end'' (e.g. transformer):
              even with an arrester installed, the travelling wave enters the
              transformer connection and causes voltage stress.
        \item The resulting stress is an overlap of the incoming surge, its
              reflection and the arrester protection rating.
        \item Equipment voltage stress therefore depends not only on the
              arrester protective level but also on the distance between
              arrester and protected device.
      \end{itemize}
    \end{column}
    \begin{column}{0.40\textwidth}
      \centering
      \topoimg[width=\linewidth,height=0.66\textheight,keepaspectratio]{\icfig{m08_image11}}
    \end{column}
  \end{columns}
\end{frame}

\begin{frame}{Energy handling capability}
  \begin{columns}[c,onlytextwidth]
    \begin{column}{0.56\textwidth}
      \begin{itemize}
        \setlength{\itemsep}{0.4em}
        \item \textbf{Simple impulse handling} --- energy injected in only a
              few \textmu s must not cause mechanical damage.
        \item \textbf{Thermal energy handling} --- maximum energy injected
              during a transient at which the arrester can still cool back to
              its nominal operating temperature.
        \item Not precisely defined in the standards, so for VSC-HVDC the
              energy loading of the arrester is an important design point.
      \end{itemize}
    \end{column}
    \begin{column}{0.40\textwidth}
      \centering
      \topoimg[width=\linewidth,height=0.66\textheight,keepaspectratio]{\icfig{m08_image12}}
    \end{column}
  \end{columns}
\end{frame}

\begin{frame}{Lightning stroke energy}
  \begin{columns}[c,onlytextwidth]
    \begin{column}{0.56\textwidth}
      \begin{itemize}
        \setlength{\itemsep}{0.35em}
        \item Energy is defined not only by the peak current but by the total
              charge $Q$ of the lightning stroke.
        \item For a 30\,kA stroke and a line surge impedance of 600\,$\Omega$,
              the travelling-wave voltage is $\sim$9000\,kV.
        \item Overvoltage causes flashover on line insulators, diverting most
              of the charge to ground before it reaches protected equipment.
        \item The most critical strokes hit substations directly --- proper
              lightning protection of equipment must be ensured.
      \end{itemize}
    \end{column}
    \begin{column}{0.40\textwidth}
      \centering
      \topoimg[width=\linewidth,height=0.66\textheight,keepaspectratio]{\icfig{m08_image13}}
    \end{column}
  \end{columns}
\end{frame}

\begin{frame}{Arrester configuration in VSC-HVDC systems}
  \begin{columns}[c,onlytextwidth]
    \begin{column}{0.48\textwidth}
      \centering
      \topoimg[width=\linewidth,height=0.66\textheight,keepaspectratio]{\icfig{m08_image15}}
    \end{column}
    \begin{column}{0.48\textwidth}
      \centering
      \topoimg[width=\linewidth,height=0.66\textheight,keepaspectratio]{\icfig{m08_image16}}
    \end{column}
  \end{columns}
  \vspace{0.3em}
  {\color{atGray}\footnotesize IEC 60071-11:2022.}
\end{frame}

\begin{frame}{Withstand voltage definition process}
  \begin{columns}[c,onlytextwidth]
    \begin{column}{0.56\textwidth}
      \begin{enumerate}
        \setlength{\itemsep}{0.4em}
        \item Determination of the representative overvoltages $U_{rp}$.
        \item Determination of the co-ordination withstand voltages
              $U_{cw}$. For HVDC it can be assumed $U_{cw} = U_{rp}$.
        \item Determination of the required withstand voltages $U_{rw}$.
        \item Determination of the specified withstand voltages $U_w$.
      \end{enumerate}
    \end{column}
    \begin{column}{0.40\textwidth}
      \centering
      \topoimg[width=\linewidth,height=0.66\textheight,keepaspectratio]{\icfig{m08_image17}}
    \end{column}
  \end{columns}
  \vspace{0.2em}
  {\color{atGray}\footnotesize IEC 60071-11:2022.}
\end{frame}

\begin{frame}{Representative overvoltages ($U_{rp}$)}
  \begin{block}{IEC 60071-11:2022}
    ``... representative overvoltages are equal to protection levels of
    arresters for directly protected equipment.''
  \end{block}
  \begin{itemize}
    \setlength{\itemsep}{0.4em}
    \item The statement is misleading: the arrester protection level is
          defined by the arrester design and the protection-coordination
          current amplitude and shape (lightning, switching).
    \item Representative overvoltages can be determined by assuming maximum
          system currents and switching impulse, or by transient-simulation
          analysis of relevant faults.
    \item Calculation and simulation results without any margin are
          considered as the representative overvoltage $U_{rp}$.
  \end{itemize}
\end{frame}

\begin{frame}{Required withstand voltages ($U_{rw}$)}
  \centering
  \topoimg[width=\linewidth,height=0.72\textheight,keepaspectratio]{\icfig{m08_image18}}
\end{frame}

\begin{frame}{Required withstand voltages ($U_{rw}$)}
  \begin{columns}[c,onlytextwidth]
    \begin{column}{0.48\textwidth}
      \centering
      \topoimg[width=\linewidth,height=0.66\textheight,keepaspectratio]{\icfig{m08_image19}}
    \end{column}
    \begin{column}{0.48\textwidth}
      \centering
      \topoimg[width=\linewidth,height=0.66\textheight,keepaspectratio]{\icfig{m08_image20}}
    \end{column}
  \end{columns}
  \vspace{0.3em}
  {\color{atGray}\footnotesize IEC 60071-11:2022.}
\end{frame}

\begin{frame}{Specified withstand voltages ($U_w$)}
  \begin{columns}[c,onlytextwidth]
    \begin{column}{0.54\textwidth}
      \begin{itemize}
        \setlength{\itemsep}{0.4em}
        \item Values are equal to or higher than $U_{rw}$.
        \item For AC equipment they correspond to the standard values in
              IEC 60071-1.
        \item For HVDC equipment they are rounded up to convenient values
              (in most cases IEC 60071-1).
      \end{itemize}
    \end{column}
    \begin{column}{0.42\textwidth}
      \centering
      \renewcommand{\arraystretch}{1.15}
      \textbf{Proposed $U_w$ / kV}\\[0.3em]
      \begin{tabular}{@{}rr@{}}
        \hline
        550  & 1300 \\
        650  & 1425 \\
        750  & 1550 \\
        850  & 1675 \\
        950  & 1800 \\
        1050 & 1950 \\
        1175 &      \\
        \hline
      \end{tabular}
    \end{column}
  \end{columns}
\end{frame}

\begin{frame}{Example for insulation point coordination}
  \begin{columns}[T,onlytextwidth]
    \begin{column}{0.48\textwidth}
      \renewcommand{\arraystretch}{1.3}
      \begin{tabular}{@{}p{4.2cm}r@{}}
        \textbf{Parameter} & \textbf{Value} \\
        \hline
        CCOV                        & 515\,kV \\
        Number of parallel columns  & 8 \\
        Energy capability           & 17\,MJ \\
      \end{tabular}
    \end{column}
    \begin{column}{0.48\textwidth}
      \renewcommand{\arraystretch}{1.3}
      \begin{tabular}{@{}p{4.4cm}r@{}}
        \textbf{Parameter} & \textbf{Value} \\
        \hline
        SIPL ($U_{rp}$ for SI) & 807\,kV @ 1\,kA \\
        LIPL ($U_{rp}$ for LI) & 872\,kV @ 5\,kA \\
      \end{tabular}
    \end{column}
  \end{columns}
\end{frame}
